% !PRINTED 593
\weberpage{593}
%<<CONT
components $x$ and $y$ are absolutely not greater than $1/2$, and hence that
\[
N(\xi-\mu) = x^2 - xy + y^2 \leqq \frac{3}{4}.
\]
On the basis of this property one can, just as in the preceding section, determine the greatest common divisor of two numbers $\alpha$, $\beta$, and thereby prove that every integer in $R(\varrho)$ can be decomposed into primes of this field in only one way, if associated numbers are not regarded as essentially different.

As in the preceding section one concludes that one obtains all the primes of the field $R(\varrho)$ by decomposing all real primes, and that a real prime either splits into two conjugate prime factors in $R(\varrho)$, $p=\pi\pi'$, or is a prime in the field $R(\varrho)$ as well.

The prime 3 splits into the two factors $(1-\varrho)$, $(1-\varrho^2)$; but these are associated, and so $3=-\varrho^2(1-\varrho)^2$ is associated with the square of a prime in $R(\varrho)$. One may choose $\sqrt{-3}$ for this prime. All primes of the form $3f+1$ can, by \S.\,172, be decomposed into two conjugate factors in $R(\varrho)$, while the primes of the form $3f+2$ are never decomposable, because an expression of the form $a^2-ab+b^2$ cannot be $\equiv 2\ (\mathrm{mod}\ 3)$. Thus one obtains all the primes in $R(\varrho)$ by seeking out the real primes of the form $3f+2$ and the conjugate factors of the remaining real primes.

If $a+b\varrho$ is a complex integer whose norm $a^2-ab+b^2$ is not divisible by 3, then of the two numbers $a$, $b$ either one must be divisible by 3 and the other not, or we must have $a\equiv b$ and consequently $a^2-ab+b^2\equiv 1\ (\mathrm{mod}\ 3)$. Among the associated numbers
\[
a+b\varrho,\quad b+a\varrho^2,\quad b-(a-b)\varrho,
\]
there is therefore always one and only one in which the coefficient of $\varrho$ (or of $\varrho^2$) is divisible by 3. Let us assume that $b=3B$ is divisible by 3, and set $N(a+b\varrho)=p$; then it follows, if in addition $2a-b=A$ is set,
\[
4p = A^2 + 27B^2,
\]
in agreement with \S.\,172. If $p$ is a real prime, and $4p$ has been represented in this form, one obtains from it the complex factors of $p$
%CONT>>
